\documentclass[
    language=english,
    doctype=article,
    institution=none,
    titlestyle=book,
]{../../omnilatex}

\RequirePackage{config/document-settings}
\addbibresource{../../bib/bibliography.bib}

\begin{document}

\maketitle

\begin{center}
    \textbf{Discussion Paper --- Comments Welcome}
\end{center}

\section*{Abstract}
This paper examines the impact of carbon pricing under the European Union
Emissions Trading System (EU ETS) on the competitiveness of energy-intensive
industries.  Using a difference-in-differences framework with firm-level data
from 2005 to 2024, we estimate the causal effect of carbon allowance costs on
production volumes, employment, and carbon leakage indicators across seven
industrial sectors.  Our results indicate that a 10\,EUR per tonne increase in
carbon allowance prices is associated with a 2.3\% reduction in production
output and a 1.1\% decrease in employment within the most exposed sectors.
However, we find limited evidence of significant carbon leakage, suggesting
that border adjustment mechanisms and free allowance allocations have
successfully mitigated relocation pressures.

\section{Introduction}
Carbon pricing has emerged as the principal market-based instrument for
reducing greenhouse gas emissions from industrial activities.  The EU ETS,
as the world's largest cap-and-trade system, covers approximately 40\% of
European greenhouse gas emissions and imposes a carbon cost on over 10,000
installations~\cite{bibid}.

The competitiveness implications of carbon pricing remain a central concern in
climate policy debates.  If carbon costs induce firms to relocate production to
jurisdictions without comparable carbon pricing---a phenomenon known as carbon
leakage---the environmental effectiveness of the policy is undermined and
domestic employment is lost~\cite{einstein1905}.  Existing theoretical models
of carbon leakage rely on assumptions about firm mobility and trade elasticities
that are difficult to test empirically~\cite{goossensLaTeXCompanion1993}.

This paper contributes to the empirical literature by exploiting the phased
implementation of the EU ETS across sectors and the variation in carbon
intensity within sectors to identify the causal effect of carbon pricing on
industrial outcomes.

\section{Institutional Background}

\subsection{The EU Emissions Trading System}
The EU ETS operates in four phases:
\begin{enumerate}
    \item \textbf{Phase I (2005--2007):} Pilot phase with free allocation
    based on historical emissions.
    \item \textbf{Phase II (2008--2012):} Tightened cap aligned with Kyoto
    Protocol targets.
    \item \textbf{Phase III (2013--2020):} Shift to auctioning and
    introduction of the Market Stability Reserve.
    \item \textbf{Phase IV (2021--2030):} Enhanced ambition with a linear
    reduction factor of 4.3\% per year and the Carbon Border Adjustment
    Mechanism (CBAM).
\end{enumerate}

\subsection{Carbon Leakage Provisions}
To address competitiveness concerns, the EU ETS provides free allowances to
sectors deemed at significant risk of carbon leakage.  The allocation is
benchmarked against the performance of the top 10\% most efficient installations
in each sector.

\section{Empirical Strategy}

\subsection{Identification}
We employ a difference-in-differences (DiD) design exploiting variation in
carbon cost exposure across sectors and over time.  The treatment intensity
is measured by the sector-specific carbon cost burden, defined as
\begin{equation}
    \text{CCB}_{s,t} = \frac{p_t \times e_s}{V_{s,t}},
    \label{eq:ccb}
\end{equation}
where $p_t$ is the carbon allowance price in year $t$, $e_s$ is the
emission intensity of sector $s$ (tonnes CO$_2$ per unit of output), and
$V_{s,t}$ is the sectoral value added.

\subsection{Estimation Equation}
Our baseline specification is
\begin{equation}
    Y_{i,s,t} = \alpha_i + \gamma_t + \delta \, \text{CCB}_{s,t} + \mathbf{X}_{i,t}'\boldsymbol{\beta} + \varepsilon_{i,s,t},
    \label{eq:did}
\end{equation}
where $Y_{i,s,t}$ is the outcome of interest for firm $i$ in sector $s$ at
time $t$, $\alpha_i$ are firm fixed effects, $\gamma_t$ are year fixed
effects, and $\mathbf{X}_{i,t}$ is a vector of time-varying firm-level
controls including firm size, capital intensity, and export share.

\section{Data}

\subsection{Source and Coverage}
We construct a panel dataset combining:
\begin{itemize}
    \item Firm-level production and financial data from the Bureau van Dijk
    Orbis database.
    \item Emissions and allowance data from the EU ETS Transaction Log.
    \item Sector-level carbon intensity estimates from Eurostat.
\end{itemize}
The final sample covers 12,847 installations across seven energy-intensive
sectors: cement, steel, aluminium, lime, glass, ceramics, and paper.

\section{Results}

\subsection{Main Findings}
Table~\ref{tab:main-results} presents the estimated effects of carbon cost
burden on key industrial outcomes.

\begin{table}[htbp]
    \centering
    \caption{Difference-in-differences estimates of the effect of a 10\,EUR/t
    CO$_2$ increase in carbon cost burden on industrial outcomes.}
    \label{tab:main-results}
    \begin{tabular}{@{}lcccc@{}}
        \toprule
        \textbf{Outcome} & \textbf{Estimate} & \textbf{SE} &
            \textbf{95\% CI} & \textbf{$N$} \\
        \midrule
        Production (log)  & $-0.023^{**}$ & 0.009 & $[-0.041,\; -0.005]$ &
            128{,}470 \\
        Employment (log)  & $-0.011^{*}$  & 0.005 & $[-0.021,\; -0.001]$ &
            128{,}470 \\
        Investment (log)  & $-0.008$      & 0.012 & $[-0.032,\; 0.016]$ &
            115{,}230 \\
        Export share      & $-0.003$      & 0.004 & $[-0.011,\; 0.005]$ &
            128{,}470 \\
        Carbon leakage    & $0.005$       & 0.007 & $[-0.009,\; 0.019]$ &
            128{,}470 \\
        \bottomrule
    \end{tabular}
    \vspace{1ex}
    \\[0.5ex]
    \footnotesize{Notes: Standard errors are clustered at the firm level.
    Significance levels: $^{*}p < 0.10$, $^{**}p < 0.05$, $^{***}p < 0.01$.}
\end{table}

The results indicate that increased carbon costs are associated with modest
reductions in production and employment, but no statistically significant
increase in carbon leakage.

\subsection{Heterogeneity Analysis}
We find substantial heterogeneity across sectors.  The cement and steel
sectors, which face the highest carbon cost burden, exhibit the largest
production responses.  In contrast, the glass and ceramics sectors show
negligible effects, likely due to their lower emission intensity and the
availability of fuel-switching options.

\section{Policy Implications}
Our findings suggest that the EU ETS competitiveness provisions---particularly
free allocation and CBAM---have been effective in preventing large-scale carbon
leakage.  However, the observed production and employment reductions in the
most exposed sectors raise concerns about the distributional consequences of
carbon pricing.

\begin{description}
    \item[Border adjustment] The phased introduction of CBAM should be
    accelerated to provide a level playing field for EU producers competing
    against imports from jurisdictions without carbon pricing.
    \item[Revenue recycling] A portion of auction revenue should be directed
    toward worker retraining and regional transition funds to address the
    employment effects documented in this study.
    \item[Free allocation reform] Free allowances should be phased out in line
    with the CBAM rollout to maintain environmental integrity while protecting
    competitiveness.
\end{description}

\section{Conclusion}
This paper provides evidence that carbon pricing under the EU ETS has modest
but significant effects on production and employment in energy-intensive
industries, without inducing large-scale carbon leakage.  These findings
support the continued use of carbon pricing as a central climate policy
instrument, provided that complementary measures are in place to address
competitiveness and distributional concerns.

Future research should examine the long-run effects of carbon pricing on
innovation and technology adoption, as well as the interaction between the EU
ETS and emerging carbon pricing systems in other jurisdictions.

\printbibliography

\end{document}
